% Source: https://artofproblemsolving.com/community/c3h2512150_1000_upvotes_received_celebration
% Stable identity: aops-topic-2512150-q16
Suppose the bird is with the Pittsburgh driver and the drivers are $D$ miles apart.
\begin{itemize}
\item \textbf{Towards Arlington.} The bird and the Arlington driver close the gap at $70+50=120$ mph, so this leg takes $\dfrac{D}{120}$ hours and the bird flies $\dfrac{70D}{120}=\dfrac{7D}{12}$ miles. Meanwhile the drivers close $100\cdot\dfrac{D}{120}$ miles, so they are now $\dfrac D6$ apart.
\item \textbf{Back towards Pittsburgh.} The closing speed is $60+50=110$ mph, so this leg takes $\dfrac{D}{660}$ hours and the bird flies $\dfrac{60D}{660}=\dfrac{D}{11}$ miles. The drivers are then $\dfrac D6-\dfrac{100D}{660}=\dfrac{D}{66}$ apart, and the bird is again with the Pittsburgh driver.
\end{itemize}
So each round trip, starting from a gap $D$, covers $\dfrac{7D}{12}+\dfrac{D}{11}=\dfrac{89D}{132}$ miles and shrinks the gap to $\dfrac D{66}$. Starting from $D=477$, the total distance is a geometric series:
\[\dfrac{89\cdot477}{132}\left(1+\dfrac1{66}+\dfrac1{66^2}+\cdots\right)=\dfrac{89\cdot477}{132}\cdot\dfrac{66}{65}=\dfrac{89\cdot477}{130}=\dfrac{42453}{130}\approx326.56.\]
(Check: the total time is $\dfrac{477\cdot66}{65}\left(\dfrac1{120}+\dfrac1{660}\right)=4.77$ hours, which is exactly when the drivers meet, $\dfrac{477}{100}$ hours after starting.)

To the nearest mile, the bird flies $\boxed{327}$ miles.
